% !TEX root = ../my-thesis.tex
%

% lange Zeilen in Code-Listings umbrechen, damit sie nicht rechts aus der Seite laufen
\lstset{breaklines=true, breakatwhitespace=false}

% ============================================================
\chapter{Python Benchmark-Skripte}
\label{app:scripts}
% ============================================================

Die folgenden vier Python-Skripte bilden die gemeinsame Rechengrundlage der Implementierungen.
In Nextflow, Pegasus und Merlin wurden sie direkt aufgerufen. Für StreamFlow wurden angepasste
Varianten verwendet, die Eingabe- und Ausgabepfade als Kommandozeilenargumente entgegennehmen
(siehe Anhang~\ref{app:streamflow:scripts}). In AiiDA wurde dieselbe Logik innerhalb der
WorkChain-Methoden nachimplementiert.

\section{generate\_input.py}
\label{app:scripts:generate}

\begin{lstlisting}[language=Python,
                   caption={generate\_input.py: Eingabedaten erzeugen},
                   label={lst:app:generate}]
#!/usr/bin/env python3

import random

numbers = [random.randint(1, 30) for i in range(10)]

with open("raw_input.txt", "w") as f:
    for n in numbers:
        f.write(f"{n}\n")
\end{lstlisting}

\section{preprocess.py}
\label{app:scripts:preprocess}

\begin{lstlisting}[language=Python,
                   caption={preprocess.py: Daten filtern},
                   label={lst:app:preprocess}]
#!/usr/bin/env python3

with open("raw_input.txt", "r") as f:
    values = [int(line.strip()) for line in f if line.strip()]

filtered = [x for x in values if x >= 10]

with open("prepared_input.txt", "w") as f:
    for x in filtered:
        f.write(f"{x}\n")
\end{lstlisting}

\section{compute.py}
\label{app:scripts:compute}

\begin{lstlisting}[language=Python,
                   caption={compute.py: Statistiken berechnen},
                   label={lst:app:compute}]
#!/usr/bin/env python3

with open("prepared_input.txt", "r") as f:
    values = [int(line.strip()) for line in f if line.strip()]

count = len(values)
total = sum(values)
mean = total / count if count > 0 else 0

with open("result.txt", "w") as f:
    f.write(f"count={count}\n")
    f.write(f"sum={total}\n")
    f.write(f"mean={mean:.2f}\n")
\end{lstlisting}

\section{postprocess.py}
\label{app:scripts:postprocess}

\begin{lstlisting}[language=Python,
                   caption={postprocess.py: Zusammenfassung erstellen},
                   label={lst:app:postprocess}]
#!/usr/bin/env python3

with open("result.txt", "r") as f:
    content = f.read()

with open("summary.txt", "w") as f:
    f.write("Summary of computation\n")
    f.write("----------------------\n")
    f.write(content)
\end{lstlisting}

% ============================================================
\chapter{Nextflow: Workflow-Dateien}
\label{app:nextflow}
% ============================================================

\section{main.nf}
\label{app:nextflow:mainnf}

\begin{lstlisting}[caption={main.nf: Vollständige Workflow-Datei (Nextflow DSL2)},
                   label={lst:app:nextflow:mainnf}]
nextflow.enable.dsl=2

process GENERATE_INPUT {
    output:
    path 'raw_input.txt'

    script:
    """
    python3 ${projectDir}/scripts/generate_input.py
    """
}

process PREPROCESS {
    input:
    path raw_file

    output:
    path 'prepared_input.txt'

    script:
    """
    python3 ${projectDir}/scripts/preprocess.py
    """
}

process COMPUTE {
    input:
    path prepared_file

    output:
    path 'result.txt'

    script:
    """
    python3 ${projectDir}/scripts/compute.py
    """
}

process POSTPROCESS {
    input:
    path result_file

    output:
    path 'summary.txt'

    script:
    """
    python3 ${projectDir}/scripts/postprocess.py
    """
}

workflow {
    raw_ch      = GENERATE_INPUT()
    prepared_ch = PREPROCESS(raw_ch)
    result_ch   = COMPUTE(prepared_ch)
    summary_ch  = POSTPROCESS(result_ch)
}
\end{lstlisting}

% ============================================================
\chapter{Pegasus: Workflow-Dateien}
\label{app:pegasus}
% ============================================================

\section{build\_workflow.py}
\label{app:pegasus:buildwf}

\begin{lstlisting}[language=Python,
                   caption={build\_workflow.py: Workflow-Komposition via Pegasus Python API},
                   label={lst:app:pegasus:buildwf}]
#!/usr/bin/env python3
from pathlib import Path
from Pegasus.api import *

base = Path(".").resolve()

tc = TransformationCatalog()
generate_t = Transformation(
    "generate",
    site="condorpool",
    pfn=(base / "scripts" / "generate_input.py").resolve().as_uri(),
    is_stageable=False,
    arch=Arch.X86_64,
    os_type=OS.MACOSX
)
preprocess_t = Transformation(
    "preprocess",
    site="condorpool",
    pfn=(base / "scripts" / "preprocess.py").resolve().as_uri(),
    is_stageable=False,
    arch=Arch.X86_64,
    os_type=OS.MACOSX
)
compute_t = Transformation(
    "compute",
    site="condorpool",
    pfn=(base / "scripts" / "compute.py").resolve().as_uri(),
    is_stageable=False,
    arch=Arch.X86_64,
    os_type=OS.MACOSX
)
postprocess_t = Transformation(
    "postprocess",
    site="condorpool",
    pfn=(base / "scripts" / "postprocess.py").resolve().as_uri(),
    is_stageable=False,
    arch=Arch.X86_64,
    os_type=OS.MACOSX
)
tc.add_transformations(
    generate_t,
    preprocess_t,
    compute_t,
    postprocess_t
)
tc.write("transformations.yml")

wf = Workflow("sequence-pipeline")
raw_file      = File("raw_input.txt")
prepared_file = File("prepared_input.txt")
result_file   = File("result.txt")
summary_file  = File("summary.txt")

generate_job = Job("generate")
generate_job.add_outputs(
    raw_file,
    stage_out=False,
    register_replica=False
)

preprocess_job = Job("preprocess")
preprocess_job.add_inputs(raw_file)
preprocess_job.add_outputs(
    prepared_file,
    stage_out=False,
    register_replica=False
)

compute_job = Job("compute")
compute_job.add_inputs(prepared_file)
compute_job.add_outputs(
    result_file,
    stage_out=False,
    register_replica=False
)

postprocess_job = Job("postprocess")
postprocess_job.add_inputs(result_file)
postprocess_job.add_outputs(
    summary_file,
    stage_out=True,
    register_replica=True
)

wf.add_jobs(
    generate_job,
    preprocess_job,
    compute_job,
    postprocess_job
)
wf.write("workflow.yml")
\end{lstlisting}

\section{transformations.yml}
\label{app:pegasus:transformations}

\begin{lstlisting}[caption={transformations.yml: Transformations-Katalog (automatisch erzeugt)},
                   label={lst:app:pegasus:transformations}]
x-pegasus:
  apiLang: python
  createdBy: Hp
  createdOn: 06-03-26T14:33:58Z
pegasus: 5.0.4
transformations:
- name: generate
  sites:
  - name: condorpool
    pfn: file:///Users/Hp/pegasus_pipeline_test/scripts/generate_input.py
    type: installed
    arch: x86_64
    os.type: macosx
- name: preprocess
  sites:
  - name: condorpool
    pfn: file:///Users/Hp/pegasus_pipeline_test/scripts/preprocess.py
    type: installed
    arch: x86_64
    os.type: macosx
- name: compute
  sites:
  - name: condorpool
    pfn: file:///Users/Hp/pegasus_pipeline_test/scripts/compute.py
    type: installed
    arch: x86_64
    os.type: macosx
- name: postprocess
  sites:
  - name: condorpool
    pfn: file:///Users/Hp/pegasus_pipeline_test/scripts/postprocess.py
    type: installed
    arch: x86_64
    os.type: macosx
\end{lstlisting}

\section{sites.yml}
\label{app:pegasus:sites}

\begin{lstlisting}[caption={sites.yml: Site-Katalog (manuell erstellt)},
                   label={lst:app:pegasus:sites}]
x-pegasus:
  apiLang: python
  createdBy: Hp
pegasus: 5.0.4

sites:
- name: local
  directories:
  - type: sharedScratch
    path: /Users/Hp/pegasus_pipeline_test/scratch
    sharedFileSystem: true
    fileServers:
    - url: file:///Users/Hp/pegasus_pipeline_test/scratch
      operation: all
  - type: localStorage
    path: /Users/Hp/pegasus_pipeline_test/output
    sharedFileSystem: true
    fileServers:
    - url: file:///Users/Hp/pegasus_pipeline_test/output
      operation: all

- name: condorpool
  arch: x86_64
  os.type: macosx
  directories:
  - type: sharedScratch
    path: /Users/Hp/pegasus_pipeline_test/scratch
    sharedFileSystem: true
    fileServers:
    - url: file:///Users/Hp/pegasus_pipeline_test/scratch
      operation: all
  profiles:
    condor:
      universe: vanilla
      periodic_remove: (JobStatus == 5) &&
        ((CurrentTime - EnteredCurrentStatus) > 10)
    pegasus:
      style: condor
      data.configuration: sharedfs
      gridstart: Kickstart
    env:
      PEGASUS_HOME: /opt/homebrew
\end{lstlisting}

\section{pegasus.properties}
\label{app:pegasus:properties}

\begin{lstlisting}[caption={pegasus.properties: Globale Konfiguration},
                   label={lst:app:pegasus:properties}]
pegasus.data.configuration = sharedfs
pegasus.transfer.arguments = -m 1
\end{lstlisting}

% ============================================================
\chapter{StreamFlow: Workflow-Dateien}
\label{app:streamflow}
% ============================================================

\section{workflow.cwl}
\label{app:streamflow:cwl}

\begin{lstlisting}[caption={workflow.cwl: Vollständige Workflow-Beschreibung (CWL)},
                   label={lst:app:streamflow:cwl}]
cwlVersion: v1.2
class: Workflow

inputs:
  generate_script:    File
  preprocess_script:  File
  compute_script:     File
  postprocess_script: File

outputs:
  final_summary:
    type: File
    outputSource: postprocess/summary_file

steps:
  generate_input:
    run:
      class: CommandLineTool
      baseCommand: python3
      inputs:
        script:
          type: File
          inputBinding:
            position: 1
        output_name:
          type: string
          default: raw_input.txt
          inputBinding:
            position: 2
      outputs:
        raw_file:
          type: File
          outputBinding:
            glob: raw_input.txt
    in:
      script: generate_script
    out: [raw_file]

  preprocess:
    run:
      class: CommandLineTool
      baseCommand: python3
      inputs:
        script:
          type: File
          inputBinding:
            position: 1
        raw_file:
          type: File
          inputBinding:
            position: 2
        output_name:
          type: string
          default: prepared_input.txt
          inputBinding:
            position: 3
      outputs:
        prepared_file:
          type: File
          outputBinding:
            glob: prepared_input.txt
    in:
      script:   preprocess_script
      raw_file: generate_input/raw_file
    out: [prepared_file]

  compute:
    run:
      class: CommandLineTool
      baseCommand: python3
      inputs:
        script:
          type: File
          inputBinding:
            position: 1
        prepared_file:
          type: File
          inputBinding:
            position: 2
        output_name:
          type: string
          default: result.txt
          inputBinding:
            position: 3
      outputs:
        result_file:
          type: File
          outputBinding:
            glob: result.txt
    in:
      script:        compute_script
      prepared_file: preprocess/prepared_file
    out: [result_file]

  postprocess:
    run:
      class: CommandLineTool
      baseCommand: python3
      inputs:
        script:
          type: File
          inputBinding:
            position: 1
        result_file:
          type: File
          inputBinding:
            position: 2
        output_name:
          type: string
          default: summary.txt
          inputBinding:
            position: 3
      outputs:
        summary_file:
          type: File
          outputBinding:
            glob: summary.txt
    in:
      script:      postprocess_script
      result_file: compute/result_file
    out: [summary_file]
\end{lstlisting}

\section{streamflow.yml}
\label{app:streamflow:yml}

\begin{lstlisting}[caption={streamflow.yml: Deployment-Konfiguration (StreamFlow)},
                   label={lst:app:streamflow:yml}]
version: v1.0

workflows:
  sequence-pipeline:
    type: cwl
    config:
      file: workflow.cwl
      settings: config.yml
    bindings:
      - step: /
        target:
          model: python-docker

models:
  python-docker:
    type: docker
    config:
      image: python:3.12-slim
\end{lstlisting}

\section{config.yml}
\label{app:streamflow:config}

\begin{lstlisting}[caption={config.yml: Eingabe-Konfiguration (StreamFlow)},
                   label={lst:app:streamflow:config}]
generate_script:
  class: File
  path: scripts/generate_input.py

preprocess_script:
  class: File
  path: scripts/preprocess.py

compute_script:
  class: File
  path: scripts/compute.py

postprocess_script:
  class: File
  path: scripts/postprocess.py
\end{lstlisting}

% ============================================================
\chapter{Merlin: Workflow-Dateien}
\label{app:merlin}
% ============================================================

\section{pipeline.yaml}
\label{app:merlin:yaml}

\begin{lstlisting}[caption={pipeline.yaml: Vollständige Workflow-Datei (Merlin)},
                   label={lst:app:merlin:yaml}]
description:
  name: merlin_pipeline
  description: Pipeline-Pattern mit Merlin.

study:
  - name: generate_input
    description: Eingabedaten erzeugen.
    run:
      cmd: python3 $(SPECROOT)/scripts/generate_input.py

  - name: preprocess
    description: Daten filtern.
    run:
      depends: [generate_input]
      cmd: |
        cp $(generate_input.workspace)/raw_input.txt .
        python3 $(SPECROOT)/scripts/preprocess.py

  - name: compute
    description: Statistiken berechnen.
    run:
      depends: [preprocess]
      cmd: |
        cp $(preprocess.workspace)/prepared_input.txt .
        python3 $(SPECROOT)/scripts/compute.py

  - name: postprocess
    description: Endgueltige Zusammenfassungsdatei erstellen.
    run:
      depends: [compute]
      cmd: |
        cp $(compute.workspace)/result.txt .
        python3 $(SPECROOT)/scripts/postprocess.py
\end{lstlisting}

% ============================================================
\chapter{AiiDA: Workflow-Dateien}
\label{app:aiida}
% ============================================================

\section{pipeline\_workchain.py}
\label{app:aiida:workchain}

\begin{lstlisting}[language=Python,
                   caption={pipeline\_workchain.py: Vollständige WorkChain-Implementierung (AiiDA)},
                   label={lst:app:aiida:workchain}]
#!/usr/bin/env python3

import random
from pathlib import Path

from aiida import load_profile
from aiida.engine import WorkChain, run_get_node
from aiida.orm import List, Dict, SinglefileData


class PipelineWorkChain(WorkChain):

    @classmethod
    def define(cls, spec):
        super().define(spec)

        spec.outline(
            cls.generate_input,
            cls.preprocess,
            cls.compute,
            cls.postprocess,
        )

        spec.output("raw_input",      valid_type=List)
        spec.output("prepared_input", valid_type=List)
        spec.output("result",         valid_type=Dict)
        spec.output("summary_file",   valid_type=SinglefileData)

    def generate_input(self):
        numbers = [random.randint(1, 30) for _ in range(10)]
        self.ctx.raw_input = List(list=numbers)

    def preprocess(self):
        values   = self.ctx.raw_input.get_list()
        filtered = [x for x in values if x >= 10]
        self.ctx.prepared_input = List(list=filtered)

    def compute(self):
        values = self.ctx.prepared_input.get_list()
        count  = len(values)
        total  = sum(values)
        mean   = total / count if count > 0 else 0
        self.ctx.result = Dict(
            dict={
                "count": count,
                "sum":   total,
                "mean":  mean,
            }
        )

    def postprocess(self):
        result = self.ctx.result.get_dict()
        summary_text = (
            "Summary of computation\n"
            "----------------------\n"
            f"count={result['count']}\n"
            f"sum={result['sum']}\n"
            f"mean={result['mean']:.2f}\n"
        )
        output_path = Path("summary.txt").resolve()
        output_path.write_text(summary_text)
        self.ctx.summary_file = SinglefileData(file=output_path)
        self.out("raw_input",      self.ctx.raw_input)
        self.out("prepared_input", self.ctx.prepared_input)
        self.out("result",         self.ctx.result)
        self.out("summary_file",   self.ctx.summary_file)


if __name__ == "__main__":
    load_profile()

    results, node = run_get_node(PipelineWorkChain)

    print("Workflow finished successfully.")
    print(f"Process PK:   {node.pk}")
    print(f"Process UUID: {node.uuid}")

    summary = results["summary_file"]
    print(f"Summary file PK:  {summary.pk}")
    print(f"Summary filename: {summary.filename}")

    with summary.open(mode="r") as handle:
        print()
        print(handle.read())
\end{lstlisting}

\section{submit\_pipeline.py}
\label{app:aiida:submit}

\begin{lstlisting}[language=Python,
                   caption={submit\_pipeline.py: Workflow-Starter fuer den Daemon-Lauf (AiiDA)},
                   label={lst:app:aiida:submit}]
#!/usr/bin/env python3

import time

from aiida import load_profile
from aiida.engine import submit
from aiida.orm import load_node

from pipeline_workchain import PipelineWorkChain


if __name__ == "__main__":
    load_profile()

    node = submit(PipelineWorkChain)

    print("Submitted PipelineWorkChain to daemon.")
    print(f"Process PK:   {node.pk}")
    print(f"Process UUID: {node.uuid}")

    while not node.is_terminated:
        node = load_node(node.pk)
        print(f"Current state: {node.process_state}")
        time.sleep(2)

    node = load_node(node.pk)

    print()
    print(f"Final state: {node.process_state}")
    print(f"Exit status: {node.exit_status}")

    if node.is_finished_ok:
        summary = node.outputs.summary_file
        print(f"Summary file PK:  {summary.pk}")
        print(f"Summary filename: {summary.filename}")
        with summary.open(mode="r") as handle:
            print()
            print(handle.read())
    else:
        print("Workflow did not finish successfully.")
\end{lstlisting}

% ============================================================
\chapter{StreamFlow: Angepasste Skripte}
\label{app:streamflow:scripts}
% ============================================================

Da StreamFlow die Workflow-Schritte in einem Docker-Container ausführt und Eingabe- und
Ausgabepfade über CWL-\texttt{inputBinding} als Kommandozeilenargumente übergibt, wurden
die Skripte für StreamFlow so angepasst, dass sie Pfade über \texttt{sys.argv}
entgegennehmen. Die Basislogik ist identisch mit den Skripten in Anhang~\ref{app:scripts}.

\section{generate\_input.py (StreamFlow)}
\label{app:streamflow:scripts:generate}

\begin{lstlisting}[language=Python,
                   caption={generate\_input.py: Angepasste Version fuer StreamFlow},
                   label={lst:app:sf:generate}]
#!/usr/bin/env python3

import sys
import random

output_file = sys.argv[1] if len(sys.argv) > 1 else "raw_input.txt"

numbers = [random.randint(1, 30) for i in range(10)]

with open(output_file, "w") as f:
    for n in numbers:
        f.write(f"{n}\n")
\end{lstlisting}

\section{preprocess.py (StreamFlow)}
\label{app:streamflow:scripts:preprocess}

\begin{lstlisting}[language=Python,
                   caption={preprocess.py: Angepasste Version fuer StreamFlow},
                   label={lst:app:sf:preprocess}]
#!/usr/bin/env python3

import sys

input_file  = sys.argv[1] if len(sys.argv) > 1 else "raw_input.txt"
output_file = sys.argv[2] if len(sys.argv) > 2 else "prepared_input.txt"

with open(input_file, "r") as f:
    values = [int(line.strip()) for line in f if line.strip()]

filtered = [x for x in values if x >= 10]

with open(output_file, "w") as f:
    for x in filtered:
        f.write(f"{x}\n")
\end{lstlisting}

\section{compute.py (StreamFlow)}
\label{app:streamflow:scripts:compute}

\begin{lstlisting}[language=Python,
                   caption={compute.py: Angepasste Version fuer StreamFlow},
                   label={lst:app:sf:compute}]
#!/usr/bin/env python3

import sys

input_file  = sys.argv[1] if len(sys.argv) > 1 else "prepared_input.txt"
output_file = sys.argv[2] if len(sys.argv) > 2 else "result.txt"

with open(input_file, "r") as f:
    values = [int(line.strip()) for line in f if line.strip()]

count = len(values)
total = sum(values)
mean  = total / count if count > 0 else 0

with open(output_file, "w") as f:
    f.write(f"count={count}\n")
    f.write(f"sum={total}\n")
    f.write(f"mean={mean:.2f}\n")
\end{lstlisting}

\section{postprocess.py (StreamFlow)}
\label{app:streamflow:scripts:postprocess}

\begin{lstlisting}[language=Python,
                   caption={postprocess.py: Angepasste Version fuer StreamFlow},
                   label={lst:app:sf:postprocess}]
#!/usr/bin/env python3

import sys

input_file  = sys.argv[1] if len(sys.argv) > 1 else "result.txt"
output_file = sys.argv[2] if len(sys.argv) > 2 else "summary.txt"

with open(input_file, "r") as f:
    content = f.read()

with open(output_file, "w") as f:
    f.write("Summary of computation\n")
    f.write("----------------------\n")
    f.write(content)
\end{lstlisting}
% ============================================================
\chapter{Scatter-Gather: Gemeinsame Python-Skripte}
\label{app:sg:scripts}
% ============================================================

Die folgenden sechs Python-Skripte bilden die Rechengrundlage des Scatter-Gather-Patterns.
Sie entsprechen den Skripten aus Anhang~\ref{app:scripts}, sind jedoch um \texttt{split.py}
und \texttt{aggregate.py} erweitert, erzeugen 100 statt zehn Eingabewerte und nehmen ihre
Ein- und Ausgabepfade als Kommandozeilenargumente (\texttt{sys.argv}) entgegen. In dieser
Form wurden sie in Nextflow, Pegasus und StreamFlow unverändert eingesetzt.
\section{generate\_input.py}
\label{app:sg:scripts:generate}

\begin{lstlisting}[language=Python,
                   caption={generate\_input.py: 100 Eingabewerte erzeugen},
                   label={lst:app:sg:generate}]
#!/usr/bin/env python3

import random
import sys

output_file = sys.argv[1] if len(sys.argv) > 1 else "raw_input.txt"

numbers = [random.randint(1, 30) for _ in range(100)]

with open(output_file, "w") as f:
    for n in numbers:
        f.write(f"{n}\n")
\end{lstlisting}

\section{preprocess.py}
\label{app:sg:scripts:preprocess}

\begin{lstlisting}[language=Python,
                   caption={preprocess.py: Daten filtern},
                   label={lst:app:sg:preprocess}]
#!/usr/bin/env python3

import sys

input_file = sys.argv[1] if len(sys.argv) > 1 else "raw_input.txt"
output_file = sys.argv[2] if len(sys.argv) > 2 else "prepared_input.txt"

with open(input_file, "r") as f:
    values = [int(line.strip()) for line in f if line.strip()]

filtered = [x for x in values if x >= 10]

with open(output_file, "w") as f:
    for x in filtered:
        f.write(f"{x}\n")
\end{lstlisting}

\section{split.py}
\label{app:sg:scripts:split}

\begin{lstlisting}[language=Python,
                   caption={split.py: Daten in Chunks aufteilen (Scatter)},
                   label={lst:app:sg:split}]
#!/usr/bin/env python3
import sys

input_file = sys.argv[1] if len(sys.argv) > 1 else "prepared_input.txt"
num_chunks = int(sys.argv[2]) if len(sys.argv) > 2 else 4

with open(input_file, "r") as f:
    values = [line.strip() for line in f if line.strip()]

chunks = [[] for _ in range(num_chunks)]
for i in range(len(values)):
    chunks[i % num_chunks].append(values[i])

for i in range(num_chunks):
    with open(f"chunk_{i+1}.txt", "w") as f:
        for value in chunks[i]:
            f.write(f"{value}\n")
\end{lstlisting}

\section{compute.py}
\label{app:sg:scripts:compute}

\begin{lstlisting}[language=Python,
                   caption={compute.py: Statistiken je Chunk berechnen},
                   label={lst:app:sg:compute}]
#!/usr/bin/env python3

import sys

input_file = sys.argv[1]
output_file = sys.argv[2]

with open(input_file, "r") as f:
    values = [int(line.strip()) for line in f if line.strip()]

count = len(values)
total = sum(values)
mean = total / count if count > 0 else 0

with open(output_file, "w") as f:
    f.write(f"chunk={input_file}\n")
    f.write(f"count={count}\n")
    f.write(f"sum={total}\n")
    f.write(f"mean={mean:.2f}\n")
\end{lstlisting}

\section{aggregate.py}
\label{app:sg:scripts:aggregate}

\begin{lstlisting}[language=Python,
                   caption={aggregate.py: Teilergebnisse zusammenführen (Gather)},
                   label={lst:app:sg:aggregate}]
#!/usr/bin/env python3

import sys

input_files = sys.argv[1:-1]
output_file = sys.argv[-1]

total_count = 0
total_sum = 0
chunks = 0

for input_file in input_files:
    chunks += 1
    values = {}

    with open(input_file, "r") as f:
        for line in f:
            if "=" in line:
                key, value = line.strip().split("=", 1)
                values[key] = value

    total_count += int(values.get("count", 0))
    total_sum += int(values.get("sum", 0))

total_mean = total_sum / total_count if total_count > 0 else 0

with open(output_file, "w") as f:
    f.write(f"chunks={chunks}\n")
    f.write(f"total_count={total_count}\n")
    f.write(f"total_sum={total_sum}\n")
    f.write(f"global_mean={total_mean:.2f}\n")
\end{lstlisting}

\section{postprocess.py}
\label{app:sg:scripts:postprocess}

\begin{lstlisting}[language=Python,
                   caption={postprocess.py: Zusammenfassung schreiben},
                   label={lst:app:sg:postprocess}]
#!/usr/bin/env python3

import sys

input_file = sys.argv[1] if len(sys.argv) > 1 else "aggregated_result.txt"
output_file = sys.argv[2] if len(sys.argv) > 2 else "summary.txt"

with open(input_file, "r") as f:
    content = f.read()

with open(output_file, "w") as f:
    f.write("Summary of Scatter-Gather computation\n")
    f.write("-------------------------------------\n")
    f.write(content)
\end{lstlisting}


% ============================================================
\chapter{Scatter-Gather: Nextflow}
\label{app:sg:nextflow}
% ============================================================

\section{main.nf}
\label{app:sg:nextflow:mainnf}
\begin{lstlisting}[caption={main.nf: Vollständige Scatter-Gather-Workflow-Datei (Nextflow DSL2)},
                   label={lst:app:sg:nextflow:mainnf}]
nextflow.enable.dsl=2

process GENERATE_INPUT {
    output:
    path 'raw_input.txt'

    script:
    """
    python3 ${projectDir}/scripts/generate_input.py raw_input.txt
    """
}

process PREPROCESS {
    input:
    path raw_file

    output:
    path 'prepared_input.txt'

    script:
    """
    python3 ${projectDir}/scripts/preprocess.py ${raw_file} prepared_input.txt
    """
}

process SPLIT {
    input:
    path prepared_file

    output:
    path 'chunk_*.txt'

    script:
    """
    python3 ${projectDir}/scripts/split.py ${prepared_file} 4
    """
}

process COMPUTE {
    input:
    path chunk_file

    output:
    path "result_${chunk_file.simpleName.replace('chunk_', '')}.txt"

    script:
    """
    chunk_id=\$(echo ${chunk_file.simpleName} | sed 's/chunk_//')
    python3 ${projectDir}/scripts/compute.py ${chunk_file} result_\${chunk_id}.txt
    """
}

process AGGREGATE {
    input:
    path result_files

    output:
    path 'aggregated_result.txt'

    script:
    """
    python3 ${projectDir}/scripts/aggregate.py ${result_files} aggregated_result.txt
    """
}

process POSTPROCESS {
    input:
    path aggregated_file

    output:
    path 'summary.txt'

    script:
    """
    python3 ${projectDir}/scripts/postprocess.py ${aggregated_file} summary.txt
    """
}

workflow {
    raw_ch = GENERATE_INPUT()

    prepared_ch = PREPROCESS(raw_ch)

    chunks_ch = SPLIT(prepared_ch)
        .flatten()

    results_ch = COMPUTE(chunks_ch)
        .collect()

    aggregated_ch = AGGREGATE(results_ch)

    summary_ch = POSTPROCESS(aggregated_ch)
}
\end{lstlisting}


% ============================================================
\chapter{Scatter-Gather: Pegasus}
\label{app:sg:pegasus}
% ============================================================

\section{build\_workflow.py}
\label{app:sg:pegasus:buildwf}
\begin{lstlisting}[language=Python,
                   caption={build\_workflow.py: Aufbau des Scatter-Gather-Workflows (Pegasus)},
                   label={lst:app:sg:pegasus:buildwf}]
#!/usr/bin/env python3

from pathlib import Path
from Pegasus.api import *

base = Path(".").resolve()

tc = TransformationCatalog()

generate_t = Transformation(
    "generate",
    site="condorpool",
    pfn=(base / "scripts" / "generate_input.py").resolve().as_uri(),
    is_stageable=False,
    arch=Arch.X86_64,
    os_type=OS.MACOSX
)

preprocess_t = Transformation(
    "preprocess",
    site="condorpool",
    pfn=(base / "scripts" / "preprocess.py").resolve().as_uri(),
    is_stageable=False,
    arch=Arch.X86_64,
    os_type=OS.MACOSX
)

split_t = Transformation(
    "split",
    site="condorpool",
    pfn=(base / "scripts" / "split.py").resolve().as_uri(),
    is_stageable=False,
    arch=Arch.X86_64,
    os_type=OS.MACOSX
)

compute_t = Transformation(
    "compute",
    site="condorpool",
    pfn=(base / "scripts" / "compute.py").resolve().as_uri(),
    is_stageable=False,
    arch=Arch.X86_64,
    os_type=OS.MACOSX
)

aggregate_t = Transformation(
    "aggregate",
    site="condorpool",
    pfn=(base / "scripts" / "aggregate.py").resolve().as_uri(),
    is_stageable=False,
    arch=Arch.X86_64,
    os_type=OS.MACOSX
)

postprocess_t = Transformation(
    "postprocess",
    site="condorpool",
    pfn=(base / "scripts" / "postprocess.py").resolve().as_uri(),
    is_stageable=False,
    arch=Arch.X86_64,
    os_type=OS.MACOSX
)

tc.add_transformations(
    generate_t,
    preprocess_t,
    split_t,
    compute_t,
    aggregate_t,
    postprocess_t
)

tc.write("transformations.yml")

wf = Workflow("scatter-gather")

raw_file = File("raw_input.txt")
prepared_file = File("prepared_input.txt")

chunk_1_file = File("chunk_1.txt")
chunk_2_file = File("chunk_2.txt")
chunk_3_file = File("chunk_3.txt")
chunk_4_file = File("chunk_4.txt")

result_1_file = File("result_1.txt")
result_2_file = File("result_2.txt")
result_3_file = File("result_3.txt")
result_4_file = File("result_4.txt")

aggregated_file = File("aggregated_result.txt")
summary_file = File("summary.txt")

generate_job = Job("generate")
generate_job.add_args("raw_input.txt")
generate_job.add_outputs(
    raw_file,
    stage_out=False,
    register_replica=False
)

preprocess_job = Job("preprocess")
preprocess_job.add_args("raw_input.txt", "prepared_input.txt")
preprocess_job.add_inputs(raw_file)
preprocess_job.add_outputs(
    prepared_file,
    stage_out=False,
    register_replica=False
)

split_job = Job("split")
split_job.add_args("prepared_input.txt", "4")
split_job.add_inputs(prepared_file)
split_job.add_outputs(
    chunk_1_file,
    chunk_2_file,
    chunk_3_file,
    chunk_4_file,
    stage_out=False,
    register_replica=False
)

compute_job_1 = Job("compute")
compute_job_1.add_args("chunk_1.txt", "result_1.txt")
compute_job_1.add_inputs(chunk_1_file)
compute_job_1.add_outputs(
    result_1_file,
    stage_out=False,
    register_replica=False
)

compute_job_2 = Job("compute")
compute_job_2.add_args("chunk_2.txt", "result_2.txt")
compute_job_2.add_inputs(chunk_2_file)
compute_job_2.add_outputs(
    result_2_file,
    stage_out=False,
    register_replica=False
)

compute_job_3 = Job("compute")
compute_job_3.add_args("chunk_3.txt", "result_3.txt")
compute_job_3.add_inputs(chunk_3_file)
compute_job_3.add_outputs(
    result_3_file,
    stage_out=False,
    register_replica=False
)

compute_job_4 = Job("compute")
compute_job_4.add_args("chunk_4.txt", "result_4.txt")
compute_job_4.add_inputs(chunk_4_file)
compute_job_4.add_outputs(
    result_4_file,
    stage_out=False,
    register_replica=False
)

aggregate_job = Job("aggregate")
aggregate_job.add_args(
    "result_1.txt",
    "result_2.txt",
    "result_3.txt",
    "result_4.txt",
    "aggregated_result.txt"
)
aggregate_job.add_inputs(
    result_1_file,
    result_2_file,
    result_3_file,
    result_4_file
)
aggregate_job.add_outputs(
    aggregated_file,
    stage_out=False,
    register_replica=False
)

postprocess_job = Job("postprocess")
postprocess_job.add_args("aggregated_result.txt", "summary.txt")
postprocess_job.add_inputs(aggregated_file)
postprocess_job.add_outputs(
    summary_file,
    stage_out=True,
    register_replica=True
)

wf.add_jobs(
    generate_job,
    preprocess_job,
    split_job,
    compute_job_1,
    compute_job_2,
    compute_job_3,
    compute_job_4,
    aggregate_job,
    postprocess_job
)

wf.write("workflow.yml")
\end{lstlisting}

\section{sites.yml}
\label{app:sg:pegasus:sites}

\begin{lstlisting}[caption={sites.yml: Site-Katalog mit sharedfs-Konfiguration (Pegasus)},
                   label={lst:app:sg:pegasus:sites}]
x-pegasus:
  apiLang: python
  createdBy: Hp
pegasus: 5.0.4

sites:
- name: local
  directories:
  - type: sharedScratch
    path: /Users/Hp/pegasus_scatter_gather_clean/scratch
    sharedFileSystem: true
    fileServers:
    - url: file:///Users/Hp/pegasus_scatter_gather_clean/scratch
      operation: all
  - type: localStorage
    path: /Users/Hp/pegasus_scatter_gather_clean/output
    sharedFileSystem: true
    fileServers:
    - url: file:///Users/Hp/pegasus_scatter_gather_clean/output
      operation: all

- name: condorpool
  arch: x86_64
  os.type: macosx
  directories:
  - type: sharedScratch
    path: /Users/Hp/pegasus_scatter_gather_clean/scratch
    sharedFileSystem: true
    fileServers:
    - url: file:///Users/Hp/pegasus_scatter_gather_clean/scratch
      operation: all
  profiles:
    condor:
      universe: vanilla
      periodic_remove: (JobStatus == 5) &&
        ((CurrentTime - EnteredCurrentStatus) > 10)
    pegasus:
      style: condor
      data.configuration: sharedfs
      gridstart: Kickstart
    env:
      PEGASUS_HOME: /opt/homebrew
\end{lstlisting}

\section{pegasus.properties}
\label{app:sg:pegasus:properties}

\begin{lstlisting}[caption={pegasus.properties: Globale Konfiguration (Pegasus)},
                   label={lst:app:sg:pegasus:properties}]
pegasus.catalog.site.file=sites.yml
pegasus.catalog.transformation.file=transformations.yml
pegasus.mode=development
pegasus.data.configuration=sharedfs
pegasus.transfer.arguments=-m 1
\end{lstlisting}


% ============================================================
\chapter{Scatter-Gather: StreamFlow}
\label{app:sg:streamflow}
% ============================================================

\section{workflow.cwl}
\label{app:sg:streamflow:cwl}
\begin{lstlisting}[caption={workflow.cwl: Vollständige Scatter-Gather-Beschreibung (CWL)},
                   label={lst:app:sg:streamflow:cwl}]
cwlVersion: v1.2
class: Workflow

inputs:
  generate_script: File
  preprocess_script: File
  split_script: File
  compute_script: File
  aggregate_script: File
  postprocess_script: File

outputs:
  final_summary:
    type: File
    outputSource: postprocess/summary_file

steps:
  generate_input:
    run:
      class: CommandLineTool
      baseCommand: python3
      inputs:
        script:
          type: File
          inputBinding:
            position: 1
        output_name:
          type: string
          default: raw_input.txt
          inputBinding:
            position: 2
      outputs:
        raw_file:
          type: File
          outputBinding:
            glob: raw_input.txt
    in:
      script: generate_script
    out: [raw_file]

  preprocess:
    run:
      class: CommandLineTool
      baseCommand: python3
      inputs:
        script:
          type: File
          inputBinding:
            position: 1
        raw_file:
          type: File
          inputBinding:
            position: 2
        output_name:
          type: string
          default: prepared_input.txt
          inputBinding:
            position: 3
      outputs:
        prepared_file:
          type: File
          outputBinding:
            glob: prepared_input.txt
    in:
      script: preprocess_script
      raw_file: generate_input/raw_file
    out: [prepared_file]

  split:
    run:
      class: CommandLineTool
      baseCommand: python3
      inputs:
        script:
          type: File
          inputBinding:
            position: 1
        prepared_file:
          type: File
          inputBinding:
            position: 2
        num_chunks:
          type: string
          default: "4"
          inputBinding:
            position: 3
      outputs:
        chunk_1:
          type: File
          outputBinding:
            glob: chunk_1.txt
        chunk_2:
          type: File
          outputBinding:
            glob: chunk_2.txt
        chunk_3:
          type: File
          outputBinding:
            glob: chunk_3.txt
        chunk_4:
          type: File
          outputBinding:
            glob: chunk_4.txt
    in:
      script: split_script
      prepared_file: preprocess/prepared_file
    out: [chunk_1, chunk_2, chunk_3, chunk_4]

  compute_1:
    run:
      class: CommandLineTool
      baseCommand: python3
      inputs:
        script:
          type: File
          inputBinding:
            position: 1
        chunk_file:
          type: File
          inputBinding:
            position: 2
        output_name:
          type: string
          default: result_1.txt
          inputBinding:
            position: 3
      outputs:
        result_file:
          type: File
          outputBinding:
            glob: result_1.txt
    in:
      script: compute_script
      chunk_file: split/chunk_1
    out: [result_file]

  compute_2:
    run:
      class: CommandLineTool
      baseCommand: python3
      inputs:
        script:
          type: File
          inputBinding:
            position: 1
        chunk_file:
          type: File
          inputBinding:
            position: 2
        output_name:
          type: string
          default: result_2.txt
          inputBinding:
            position: 3
      outputs:
        result_file:
          type: File
          outputBinding:
            glob: result_2.txt
    in:
      script: compute_script
      chunk_file: split/chunk_2
    out: [result_file]

  compute_3:
    run:
      class: CommandLineTool
      baseCommand: python3
      inputs:
        script:
          type: File
          inputBinding:
            position: 1
        chunk_file:
          type: File
          inputBinding:
            position: 2
        output_name:
          type: string
          default: result_3.txt
          inputBinding:
            position: 3
      outputs:
        result_file:
          type: File
          outputBinding:
            glob: result_3.txt
    in:
      script: compute_script
      chunk_file: split/chunk_3
    out: [result_file]

  compute_4:
    run:
      class: CommandLineTool
      baseCommand: python3
      inputs:
        script:
          type: File
          inputBinding:
            position: 1
        chunk_file:
          type: File
          inputBinding:
            position: 2
        output_name:
          type: string
          default: result_4.txt
          inputBinding:
            position: 3
      outputs:
        result_file:
          type: File
          outputBinding:
            glob: result_4.txt
    in:
      script: compute_script
      chunk_file: split/chunk_4
    out: [result_file]

  aggregate:
    run:
      class: CommandLineTool
      baseCommand: python3
      inputs:
        script:
          type: File
          inputBinding:
            position: 1
        result_1:
          type: File
          inputBinding:
            position: 2
        result_2:
          type: File
          inputBinding:
            position: 3
        result_3:
          type: File
          inputBinding:
            position: 4
        result_4:
          type: File
          inputBinding:
            position: 5
        output_name:
          type: string
          default: aggregated_result.txt
          inputBinding:
            position: 6
      outputs:
        aggregated_file:
          type: File
          outputBinding:
            glob: aggregated_result.txt
    in:
      script: aggregate_script
      result_1: compute_1/result_file
      result_2: compute_2/result_file
      result_3: compute_3/result_file
      result_4: compute_4/result_file
    out: [aggregated_file]

  postprocess:
    run:
      class: CommandLineTool
      baseCommand: python3
      inputs:
        script:
          type: File
          inputBinding:
            position: 1
        aggregated_file:
          type: File
          inputBinding:
            position: 2
        output_name:
          type: string
          default: summary.txt
          inputBinding:
            position: 3
      outputs:
        summary_file:
          type: File
          outputBinding:
            glob: summary.txt
    in:
      script: postprocess_script
      aggregated_file: aggregate/aggregated_file
    out: [summary_file]
\end{lstlisting}

\section{streamflow.yml}
\label{app:sg:streamflow:yml}

\begin{lstlisting}[caption={streamflow.yml: Deployment mit replicas (StreamFlow)},
                   label={lst:app:sg:streamflow:yml}]
version: v1.0

workflows:
  scatter-gather:
    type: cwl
    config:
      file: workflow.cwl
      settings: config.yml
    bindings:
      - step: /
        target:
          model: python-docker

models:
  python-docker:
    type: docker
    config:
      image: python:3.12-slim
      replicas: 4
\end{lstlisting}

\section{config.yml}
\label{app:sg:streamflow:config}

\begin{lstlisting}[caption={config.yml: Eingabe-Konfiguration (StreamFlow)},
                   label={lst:app:sg:streamflow:config}]
generate_script:
  class: File
  path: scripts/generate_input.py

preprocess_script:
  class: File
  path: scripts/preprocess.py

split_script:
  class: File
  path: scripts/split.py

compute_script:
  class: File
  path: scripts/compute.py

aggregate_script:
  class: File
  path: scripts/aggregate.py

postprocess_script:
  class: File
  path: scripts/postprocess.py
\end{lstlisting}


\chapter{Scatter-Gather: Merlin}
\label{app:sg:merlin}
% ============================================================

\section{scatter\_gather.yaml}
\label{app:sg:merlin:yaml}

\begin{lstlisting}[caption={scatter\_gather.yaml: Vollständige Workflow-Datei (Merlin, Scatter-Gather)},
                   label={lst:app:sg:merlin:yaml}]
description:
  name: merlin_scatter_gather
  description: Scatter Gather Pattern mit Merlin

study:
  - name: generate_input
    description: erzeugt Eingabedaten
    run:
      cmd: python3 $(SPECROOT)/scripts/generate_input.py raw_input.txt

  - name: preprocess
    description: Daten ausfiltern
    run:
      depends: [generate_input]
      cmd: |
        cp $(generate_input.workspace)/raw_input.txt .
        python3 $(SPECROOT)/scripts/preprocess.py raw_input.txt prepared_input.txt

  - name: split
    description: verarbeitete Daten in vier Chunks aufteilen
    run:
      depends: [preprocess]
      cmd: |
        cp $(preprocess.workspace)/prepared_input.txt .
        python3 $(SPECROOT)/scripts/split.py prepared_input.txt 4

  - name: compute_1
    description: Erster Chunk verarbeiten
    run:
      depends: [split]
      cmd: |
        cp $(split.workspace)/chunk_1.txt .
        python3 $(SPECROOT)/scripts/compute.py chunk_1.txt result_1.txt

  - name: compute_2
    description: Zweiter Chunk verarbeiten
    run:
      depends: [split]
      cmd: |
        cp $(split.workspace)/chunk_2.txt .
        python3 $(SPECROOT)/scripts/compute.py chunk_2.txt result_2.txt

  - name: compute_3
    description: Dritter Chunk verarbeiten
    run:
      depends: [split]
      cmd: |
        cp $(split.workspace)/chunk_3.txt .
        python3 $(SPECROOT)/scripts/compute.py chunk_3.txt result_3.txt

  - name: compute_4
    description: Vierter Chunk verarbeiten
    run:
      depends: [split]
      cmd: |
        cp $(split.workspace)/chunk_4.txt .
        python3 $(SPECROOT)/scripts/compute.py chunk_4.txt result_4.txt

  - name: aggregate
    description: Ergebnisse der vier Compute-Schritte zusammenführen.
    run:
      depends: [compute_1, compute_2, compute_3, compute_4]
      cmd: |
        cp $(compute_1.workspace)/result_1.txt .
        cp $(compute_2.workspace)/result_2.txt .
        cp $(compute_3.workspace)/result_3.txt .
        cp $(compute_4.workspace)/result_4.txt .
        python3 $(SPECROOT)/scripts/aggregate.py result_1.txt result_2.txt result_3.txt result_4.txt aggregated_result.txt

  - name: postprocess
    description: Finale Zusammenfassungsdatei erstellen.
    run:
      depends: [aggregate]
      cmd: |
        cp $(aggregate.workspace)/aggregated_result.txt .
        python3 $(SPECROOT)/scripts/postprocess.py aggregated_result.txt summary.txt
\end{lstlisting}

% ============================================================
\chapter{Scatter-Gather: AiiDA}
\label{app:sg:aiida}
% ============================================================

\section{scatter\_gather\_workchain.py}
\label{app:sg:aiida:workchain}

\begin{lstlisting}[language=Python,
                   caption={scatter\_gather\_workchain.py: Vollständige WorkChain-Implementierung (AiiDA, Scatter-Gather)},
                   label={lst:app:sg:aiida:workchain}]

import random
import time
import tempfile
from pathlib import Path

from aiida.engine import WorkChain, ToContext, calcfunction
from aiida.orm import Int, Float, List, Dict, SinglefileData


@calcfunction
def generate_input(count):
    numbers = [random.randint(1, 30) for _ in range(count.value)]
    return List(list=numbers)


@calcfunction
def preprocess(numbers):
    values = numbers.get_list()
    filtered = [x for x in values if x >= 10]
    return List(list=filtered)


@calcfunction
def split_data(numbers, num_chunks):
    values = numbers.get_list()
    n = num_chunks.value

    chunks = [[] for _ in range(n)]

    for i, value in enumerate(values):
        chunks[i % n].append(value)

    return List(list=chunks)


@calcfunction
def compute_chunk(chunk, index, sleep):
    """Verarbeitet Chunk.
    sleep für den Nachweis der Parallelitä
    """
    if sleep.value > 0:
        time.sleep(sleep.value)

    values = chunk.get_list()

    count = len(values)
    total = sum(values)
    mean = total / count if count > 0 else 0

    return Dict(
        dict={
            "chunk": f"chunk_{index.value}.txt",
            "count": count,
            "sum": total,
            "mean": round(mean, 2),
        }
    )


@calcfunction
def aggregate(**results):
    total_count = 0
    total_sum = 0
    chunks = 0

    for label in sorted(results.keys()):
        data = results[label].get_dict()
        chunks += 1
        total_count += data["count"]
        total_sum += data["sum"]

    global_mean = total_sum / total_count if total_count > 0 else 0

    return Dict(
        dict={
            "chunks": chunks,
            "total_count": total_count,
            "total_sum": total_sum,
            "global_mean": round(global_mean, 2),
        }
    )


@calcfunction
def postprocess(result):
    r = result.get_dict()

    text = (
        "Summary of Scatter-Gather computation\n"
        "-------------------------------------\n"
        f"chunks={r['chunks']}\n"
        f"total_count={r['total_count']}\n"
        f"total_sum={r['total_sum']}\n"
        f"global_mean={r['global_mean']:.2f}\n"
    )

    with tempfile.NamedTemporaryFile("w", suffix=".txt", delete=False) as handle:
        handle.write(text)
        path = handle.name

    node = SinglefileData(file=path, filename="summary.txt")
    Path(path).unlink(missing_ok=True)

    return node


class ComputeChunkWorkChain(WorkChain):
    """Verarbeitet genau einen Chunk als eigener AiiDA-Prozess."""

    @classmethod
    def define(cls, spec):
        super().define(spec)

        spec.input("chunk", valid_type=List)
        spec.input("index", valid_type=Int)
        spec.input("sleep", valid_type=Float, default=lambda: Float(0.0))

        spec.outline(cls.run_compute)

        spec.output("result", valid_type=Dict)

    def run_compute(self):
        index = self.inputs.index.value
        sleep = self.inputs.sleep.value

        self.report(f"START compute_{index}: sleep={sleep}s")

        result = compute_chunk(
            self.inputs.chunk,
            self.inputs.index,
            self.inputs.sleep,
        )

        self.report(f"END compute_{index}")

        self.out("result", result)


class ScatterGatherWorkChain(WorkChain):
    """generate -> preprocess -> split -> compute parallel -> aggregate -> postprocess."""

    @classmethod
    def define(cls, spec):
        super().define(spec)

        spec.input("count", valid_type=Int, default=lambda: Int(100))
        spec.input("num_chunks", valid_type=Int, default=lambda: Int(4))
        spec.input("compute_sleep", valid_type=Float, default=lambda: Float(5.0))

        spec.outline(
            cls.generate,
            cls.split,
            cls.submit_compute,
            cls.gather,
        )

        spec.output("raw_input", valid_type=List)
        spec.output("prepared_input", valid_type=List)
        spec.output("chunks", valid_type=List)
        spec.output("aggregated_result", valid_type=Dict)
        spec.output("summary_file", valid_type=SinglefileData)

        spec.exit_code(
            400,
            "ERROR_COMPUTE_FAILED",
            message="Ein Compute-Schritt ist fehlgeschlagen.",
        )

    def generate(self):
        self.ctx.raw_input = generate_input(self.inputs.count)
        self.ctx.prepared_input = preprocess(self.ctx.raw_input)

    def split(self):
        self.ctx.chunks = split_data(
            self.ctx.prepared_input,
            self.inputs.num_chunks,
        )

    def submit_compute(self):
        chunks = self.ctx.chunks.get_list()
        futures = {}

        for i, chunk in enumerate(chunks, start=1):
            future = self.submit(
                ComputeChunkWorkChain,
                chunk=List(list=chunk),
                index=Int(i),
                sleep=self.inputs.compute_sleep,
            )

            futures[f"compute_{i}"] = future
            self.report(f"submitted compute_{i} <{future.pk}> fuer chunk {i}")

        return ToContext(**futures)

    def gather(self):
        num_chunks = self.inputs.num_chunks.value
        results = {}

        for i in range(1, num_chunks + 1):
            child = self.ctx[f"compute_{i}"]

            if not child.is_finished_ok:
                self.report(f"compute_{i} <{child.pk}> fehlgeschlagen")
                return self.exit_codes.ERROR_COMPUTE_FAILED

            results[f"compute_{i}"] = child.outputs.result

        aggregated = aggregate(**results)
        summary_file = postprocess(aggregated)

        self.out("raw_input", self.ctx.raw_input)
        self.out("prepared_input", self.ctx.prepared_input)
        self.out("chunks", self.ctx.chunks)
        self.out("aggregated_result", aggregated)
        self.out("summary_file", summary_file)
\end{lstlisting}

\section{submit\_scatter\_gather.py}
\label{app:sg:aiida:submit}

\begin{lstlisting}[language=Python,
                   caption={submit\_scatter\_gather.py: Submit-Skript fuer den Daemon-Lauf (AiiDA, Scatter-Gather)},
                   label={lst:app:sg:aiida:submit}]
"""Startet die parallele ScatterGatherWorkChain über den AiiDA-Daemon."""

import time

from aiida import load_profile
from aiida.engine import submit
from aiida.orm import Int, Float, load_node

from scatter_gather_workchain import ScatterGatherWorkChain


load_profile()

start = time.time()

node = submit(
    ScatterGatherWorkChain,
    count=Int(100),
    num_chunks=Int(4),
    compute_sleep=Float(5.0),
)

print(f"Submitted ScatterGatherWorkChain to daemon.")
print(f"Process PK: {node.pk}")
print(f"Process UUID: {node.uuid}")
print("\nWaiting for workflow to finish...")

while not node.is_terminated:
    node = load_node(node.pk)
    print(f"Current state: {node.process_state}")
    time.sleep(1)

node = load_node(node.pk)
elapsed = time.time() - start

print("\nWorkflow finished.")
print(f"Process state: {node.process_state}")
print(f"Exit status: {node.exit_status}")
print(f"Gesamtzeit: {elapsed:.1f}s")
print(f"Pk {node.pk}")

if node.is_finished_ok:
    print("\naggregated_result:")
    print(node.outputs.aggregated_result.get_dict())

    print("\nsummary.txt:")
    print(node.outputs.summary_file.get_content())
else:
    print("\nWorkflow nicht erfolgreich beendet.")
    print(f"Details: verdi process report {node.pk}")
\end{lstlisting}
% ============================================================
\chapter{Laufzeit-Logs zum Nachweis der Parallelität}
\label{app:logs}
% ============================================================

Dieser Anhang enthält die Log-Ausschnitte. Sie belegen für Pegasus, Merlin und AiiDA, dass sich die vier Compute-Schritte zeitlich
überlappend, also parallel liefen. Für Nextflow und StreamFlow ist der Nachweis über die HTML-Timeline, welche im Anhang~\ref{app:figures} abgebildet sind.

\section{Pegasus: Laufzeiten aus den .out-Dateien}
\label{app:logs:pegasus}

\begin{lstlisting}[caption={Start- und Endzeiten der vier compute-Jobs (Pegasus)},label={lst:logs:pegasus}]
compute_ID0000006   START 14:55:59.289   END 14:56:04.296
compute_ID0000005   START 14:56:00.081   END 14:56:05.084
compute_ID0000004   START 14:56:00.354   END 14:56:05.360
compute_ID0000007   START 14:56:01.109   END 14:56:06.123
\end{lstlisting}

Alle vier Jobs warteten künstlich 5 Sekunden. Bei serieller Ausführung wären das rund
20 Sekunden gewesen; tatsächlich lagen erster Start (14:55:59) und letztes Ende (14:56:06) nur
etwa 7 Sekunden auseinander.

\section{Merlin: Worker-Logs}
\label{app:logs:merlin}

\begin{lstlisting}[caption={Start- und Endzeiten der vier compute-Schritte (Merlin)},label={lst:logs:merlin}]
compute_4   Start 21:01:44,163
compute_1   Start 21:01:44,167
compute_3   Start 21:01:44,170
compute_2   Start 21:01:44,174

Ende der vier Schritte zwischen 21:01:49,499 und 21:01:50,324
aggregate   Start 21:01:50,365
\end{lstlisting}

Die vier Schritte starteten innerhalb weniger Millisekunden gemeinsam. Der
\texttt{aggregate}-Schritt begann erst, nachdem alle vier beendet waren.

\section{AiiDA: verdi process report (PK 243)}
\label{app:logs:aiida}

\begin{lstlisting}[caption={Auszug aus verdi process report fuer ScatterGatherWorkChain PK 243 (AiiDA)},label={lst:logs:aiida}]
2026-06-15 00:46:38 [243|ScatterGatherWorkChain|submit_compute]: submitted compute_1 <252> fuer chunk 1
2026-06-15 00:46:38 [243|ScatterGatherWorkChain|submit_compute]: submitted compute_2 <255> fuer chunk 2
2026-06-15 00:46:39 [243|ScatterGatherWorkChain|submit_compute]: submitted compute_3 <258> fuer chunk 3
2026-06-15 00:46:39 [243|ScatterGatherWorkChain|submit_compute]: submitted compute_4 <262> fuer chunk 4
2026-06-15 00:46:39   [252|ComputeChunkWorkChain|run_compute]: START compute_1: sleep=5.0s
2026-06-15 00:46:39   [255|ComputeChunkWorkChain|run_compute]: START compute_2: sleep=5.0s
2026-06-15 00:46:39   [258|ComputeChunkWorkChain|run_compute]: START compute_3: sleep=5.0s
2026-06-15 00:46:39   [262|ComputeChunkWorkChain|run_compute]: START compute_4: sleep=5.0s
2026-06-15 00:46:44   [252|ComputeChunkWorkChain|run_compute]: END compute_1
2026-06-15 00:46:44   [255|ComputeChunkWorkChain|run_compute]: END compute_2
2026-06-15 00:46:44   [258|ComputeChunkWorkChain|run_compute]: END compute_3
2026-06-15 00:46:44   [262|ComputeChunkWorkChain|run_compute]: END compute_4
\end{lstlisting}

Die vier Child-Prozesse (PK 252, 255, 258 und 262) starteten zwischen 00:46:38 und 00:46:39 und
endeten gemeinsam um 00:46:44. Die im Submit-Skript gemessene Gesamtzeit lag bei etwa
8,7 Sekunden.

% ============================================================
\chapter{Abbildungen: Reports, Timelines und DAG}
\label{app:figures}
% ============================================================

Dieser Anhang zeigt die relevanten Bereiche der erzeugten HTML-Reports und Timelines sowie den
Pegasus-Ausführungsgraphen. 

\begin{figure}[ht]
  \centering
  \includegraphics[width=0.9\textwidth]{content/figs/nextflow_timeline.png}
  \caption{Nextflow-Timeline: Die vier COMPUTE-Aufrufe überlappen sich zeitlich.}
  \label{fig:nextflow:timeline}
\end{figure}

\begin{figure}[ht]
  \centering
  \includegraphics[width=0.9\textwidth]{content/figs/streamflow_timeline.png}
  \caption{StreamFlow-Report mit \texttt{replicas: 4}: Die vier compute-Schritte laufen zeitlich überlappend.}
  \label{fig:streamflow:report}
\end{figure}

\begin{figure}[ht]
  \centering
  \includegraphics[width=0.7\textwidth]{content/figs/pegasus_dag.png}
  \caption{Pegasus-DAG: Die vier compute-Jobs liegen als unabhängige Zweige zwischen \texttt{split} und \texttt{aggregate}.}
  \label{fig:pegasus:dag}
\end{figure}